% !TeX program = pdflatex
\documentclass[11pt,a4paper]{article}

\usepackage[T1]{fontenc}
\usepackage[utf8]{inputenc}
\usepackage[margin=2.5cm]{geometry}
\usepackage{amsmath,amssymb}
\usepackage{booktabs}
\usepackage{array}
\usepackage{xcolor}
\usepackage{listings}
\usepackage[hidelinks]{hyperref}

%%%%%%%%%%%%%%%%%%%%%%%%%%%%%%%%%%%%%%%%%%%%%%%%%%%%%%%%%%%%%%%%%%%%%%%%%%%%%%%
% Listing style for Turtle / OWL functional syntax
%%%%%%%%%%%%%%%%%%%%%%%%%%%%%%%%%%%%%%%%%%%%%%%%%%%%%%%%%%%%%%%%%%%%%%%%%%%%%%%
\definecolor{ttlkeyword}{rgb}{0.10,0.30,0.65}
\definecolor{ttlcomment}{rgb}{0.45,0.45,0.45}
\definecolor{ttlstring}{rgb}{0.60,0.15,0.15}
\definecolor{ttlbg}{gray}{0.97}

\lstdefinelanguage{turtle}{
  morekeywords={a,rdf:type,rdfs:subClassOf,rdfs:subPropertyOf,rdfs:domain,
    rdfs:range,rdfs:Datatype,owl:Class,owl:Restriction,owl:Thing,owl:Nothing,
    owl:onProperty,owl:someValuesFrom,owl:allValuesFrom,owl:hasValue,owl:hasSelf,
    owl:intersectionOf,owl:unionOf,owl:complementOf,owl:oneOf,
    owl:equivalentClass,owl:disjointWith,owl:inverseOf,
    owl:propertyChainAxiom,owl:propertyDisjointWith,owl:hasKey,
    owl:TransitiveProperty,owl:SymmetricProperty,owl:AsymmetricProperty,
    owl:ReflexiveProperty,owl:IrreflexiveProperty,
    owl:FunctionalProperty,owl:InverseFunctionalProperty,
    owl:ObjectProperty,owl:DatatypeProperty,owl:NamedIndividual,
    owl:minCardinality,owl:maxCardinality,owl:cardinality,
    owl:minQualifiedCardinality,owl:maxQualifiedCardinality,
    owl:qualifiedCardinality,owl:onClass,owl:onDataRange,owl:onDatatype,
    owl:withRestrictions,xsd:integer,xsd:string,xsd:boolean,
    xsd:nonNegativeInteger,xsd:minInclusive,xsd:maxInclusive},
  sensitive=true,
  morecomment=[l]{\#},
  morestring=[b]",
}

\lstset{
  language=turtle,
  basicstyle=\ttfamily\footnotesize,
  keywordstyle=\color{ttlkeyword},
  commentstyle=\color{ttlcomment}\itshape,
  stringstyle=\color{ttlstring},
  backgroundcolor=\color{ttlbg},
  frame=leftline,
  framesep=6pt,
  xleftmargin=10pt,
  rulecolor=\color{gray},
  breaklines=true,
  columns=fullflexible,
  keepspaces=true,
  showstringspaces=false,
  aboveskip=6pt,
  belowskip=6pt,
}

%%%%%%%%%%%%%%%%%%%%%%%%%%%%%%%%%%%%%%%%%%%%%%%%%%%%%%%%%%%%%%%%%%%%%%%%%%%%%%%
% Shorthands for DL notation
%%%%%%%%%%%%%%%%%%%%%%%%%%%%%%%%%%%%%%%%%%%%%%%%%%%%%%%%%%%%%%%%%%%%%%%%%%%%%%%
\newcommand{\dl}[1]{\ensuremath{\mathcal{#1}}}
\newcommand{\ALC}{\dl{ALC}}
\newcommand{\EL}{\dl{EL}}
\newcommand{\AL}{\dl{AL}}
\newcommand{\SROIQ}{\dl{SROIQ}}
\newcommand{\SHOIN}{\dl{SHOIN}}
\newcommand{\SHIF}{\dl{SHIF}}
\newcommand{\dllite}[1]{\ensuremath{\textit{DL-Lite}_{#1}}}
\newcommand{\dlliteplain}{\textit{DL-Lite}}
\newcommand{\cn}[1]{\mathsf{#1}}      % concept name
\newcommand{\rn}[1]{\mathsf{#1}}      % role name
\newcommand{\ind}[1]{\mathsf{#1}}     % individual name
\newcommand{\sub}{\sqsubseteq}
\newcommand{\eqv}{\equiv}
\newcommand{\inv}[1]{#1^{-}}

% Environment pairing a DL axiom with its OWL serialisation
\newcommand{\dlaxiom}[1]{%
  \begin{displaymath}#1\end{displaymath}}

\title{Expressivity of Description Logics\\and the Corresponding OWL Profiles}
\author{}
\date{}

\begin{document}
\maketitle
\vspace{-2\baselineskip}

\begin{abstract}
\noindent This document gives a letter-by-letter account of Description Logic (DL)
naming conventions. For each label it provides a description of the constructor,
an example axiom in DL syntax, and the corresponding OWL~2 encoding in Turtle.
It closes with the mapping between named DLs and OWL species, with a
comparison of the three OWL~2 profiles, with an analysis of the rule-based
entailment regimes offered by GraphDB --- including an empirical check of what
each ruleset actually infers --- and with a survey of openly available
reasoners.
\end{abstract}

\begin{center}
\footnotesize
\textbf{DL identifiers at a glance}\\[4pt]
\begin{tabular}{@{}clll@{}}
\toprule
\textbf{Id} & \textbf{Construct} & \textbf{DL syntax} & \textbf{OWL~2 vocabulary}\\
\midrule
\multicolumn{4}{@{}l}{\emph{Base languages}}\\
\AL{} & attributive language & $\top, \bot, \neg A, C \sqcap D, \forall R.C, \exists R.\top$
  & \texttt{intersectionOf}, \texttt{allValuesFrom}, \dots\\
\EL{} & existential language & $\top, C \sqcap D, \exists R.C$
  & \texttt{intersectionOf}, \texttt{someValuesFrom}\\
\dl{FL}$_0$ & frame language & $C \sqcap D, \forall R.C$
  & \texttt{intersectionOf}, \texttt{allValuesFrom}\\
\dl{S} & \ALC{} + transitive roles & $\ALC$, $\mathsf{Trans}(R)$
  & \texttt{TransitiveProperty}\\
\midrule
\multicolumn{4}{@{}l}{\emph{Constructor letters}}\\
\dl{C} & complex negation & $\neg C$ & \texttt{complementOf}\\
\dl{U} & union & $C \sqcup D$ & \texttt{unionOf}\\
\dl{E} & full existential & $\exists R.C$ & \texttt{someValuesFrom}\\
\dl{H} & role hierarchy & $R \sub S$ & \texttt{subPropertyOf}\\
\dl{R} & complex role inclusions & $R_1 \circ R_2 \sub S$, $\exists R.\mathsf{Self}$
  & \texttt{propertyChainAxiom}, \texttt{hasSelf}\\
\dl{O} & nominals & $\{a, b\}$ & \texttt{oneOf}, \texttt{hasValue}\\
\dl{I} & inverse roles & $\inv{R}$ & \texttt{inverseOf}\\
\dl{N} & unqualified cardinality & ${\geq}n\,R$, ${\leq}n\,R$
  & \texttt{min/maxCardinality}\\
\dl{Q} & qualified cardinality & ${\geq}n\,R.C$, ${\leq}n\,R.C$
  & \texttt{min/maxQualifiedCardinality}\\
\dl{F} & functional roles & $\top \sub {\leq}1\,R$
  & \texttt{FunctionalProperty}\\
$(\mathcal{D})$ & datatypes & $\exists U.d$, facets
  & \texttt{DatatypeProperty}, \texttt{withRestrictions}\\
\midrule
\multicolumn{4}{@{}l}{\emph{Families}}\\
\dlliteplain{} & FOL-rewritable fragments & $B \sub C$, $B ::= A \mid \exists R$
  & OWL~2~QL ($\dllite{\mathcal{R}}$)\\
\bottomrule
\end{tabular}
\end{center}

\newpage
\tableofcontents
\bigskip

%%%%%%%%%%%%%%%%%%%%%%%%%%%%%%%%%%%%%%%%%%%%%%%%%%%%%%%%%%%%%%%%%%%%%%%%%%%%%%%
\section{How the naming works}

A DL name is built by concatenating letters: a \emph{base language}
(\AL{}, \EL{}, \dl{FL}, or the shorthand \dl{S}) followed by one letter per
additional constructor. Thus
\[
  \SROIQ^{(\mathcal{D})} \;=\;
  \underbrace{\dl{S}}_{\ALC + \text{trans.}} +
  \underbrace{\dl{R}}_{\text{role axioms}} +
  \underbrace{\dl{O}}_{\text{nominals}} +
  \underbrace{\dl{I}}_{\text{inverses}} +
  \underbrace{\dl{Q}}_{\text{qual.\ card.}} +
  \underbrace{(\mathcal{D})}_{\text{datatypes}} .
\]
Throughout, $A$ denotes an atomic concept (class name), $C, D$ arbitrary concept
expressions, $R, S$ roles (properties), and $a, b$ individuals.
All Turtle snippets assume the usual \texttt{rdf:}, \texttt{rdfs:},
\texttt{owl:}, \texttt{xsd:} prefixes and a local prefix \texttt{:}.

%%%%%%%%%%%%%%%%%%%%%%%%%%%%%%%%%%%%%%%%%%%%%%%%%%%%%%%%%%%%%%%%%%%%%%%%%%%%%%%
\section{Base languages}

\subsection*{\dl{FL}$_0$ / \dl{FL}$^-$ --- frame-based languages}
\dl{FL}$_0$ offers conjunction and value restriction only; \dl{FL}$^-$ adds the
limited existential $\exists R.\top$. Both are of mainly historical interest and
neither corresponds to a named OWL fragment.
\dlaxiom{\cn{Parent} \sub \cn{Person} \sqcap \forall \rn{hasChild}.\cn{Person}}

\subsection*{\AL{} --- Attributive Language}
The classical baseline: atomic concepts, $\top$, $\bot$, \emph{atomic} negation
(only $\neg A$ for a class name $A$), conjunction, value restriction
$\forall R.C$, and \emph{unqualified} existential $\exists R.\top$.
\dlaxiom{\cn{Parent} \sub \cn{Person} \sqcap \neg\cn{Robot} \sqcap
  \forall \rn{hasChild}.\cn{Person} \sqcap \exists \rn{hasChild}.\top}
\begin{lstlisting}
:Parent rdfs:subClassOf [
  a owl:Class ;
  owl:intersectionOf (
    :Person
    [ a owl:Class ; owl:complementOf :Robot ]
    [ a owl:Restriction ; owl:onProperty :hasChild ;
      owl:allValuesFrom :Person ]
    [ a owl:Restriction ; owl:onProperty :hasChild ;
      owl:someValuesFrom owl:Thing ]
  )
] .
\end{lstlisting}

\subsection*{\EL{} --- Existential Language}
$\top$, conjunction, and \emph{full} existential restriction $\exists R.C$.
No negation, no disjunction, no value restriction. This is the base of OWL~2~EL.
\dlaxiom{\cn{Parent} \eqv \cn{Person} \sqcap \exists \rn{hasChild}.\cn{Person}}
\begin{lstlisting}
:Parent owl:equivalentClass [
  a owl:Class ;
  owl:intersectionOf (
    :Person
    [ a owl:Restriction ; owl:onProperty :hasChild ;
      owl:someValuesFrom :Person ]
  )
] .
\end{lstlisting}

%%%%%%%%%%%%%%%%%%%%%%%%%%%%%%%%%%%%%%%%%%%%%%%%%%%%%%%%%%%%%%%%%%%%%%%%%%%%%%%
\section{Constructor letters}

\subsection*{\dl{C} --- complex (full) concept negation}
Negation of arbitrary concept expressions, not just of class names.
$\AL + \dl{C} = \ALC$, which is closed under the Boolean operators; hence
\dl{U} and \dl{E} come for free.
\dlaxiom{\cn{Childless} \eqv \neg \exists \rn{hasChild}.\top}
\begin{lstlisting}
:Childless owl:equivalentClass [
  a owl:Class ;
  owl:complementOf [ a owl:Restriction ;
                     owl:onProperty :hasChild ;
                     owl:someValuesFrom owl:Thing ]
] .
\end{lstlisting}

\subsection*{\dl{U} --- concept union}
Disjunction of concept expressions. Redundant once \dl{C} is available
(De Morgan), which is why it rarely appears in modern names.
\dlaxiom{\cn{Person} \eqv \cn{Man} \sqcup \cn{Woman}}
\begin{lstlisting}
:Person owl:equivalentClass [
  a owl:Class ; owl:unionOf ( :Man :Woman )
] .
\end{lstlisting}

\subsection*{\dl{E} --- full existential restriction}
$\exists R.C$ with an arbitrary filler, as opposed to \AL{}'s
$\exists R.\top$. Also redundant given \dl{C} and $\forall$.
\dlaxiom{\cn{Researcher} \sub \exists \rn{worksFor}.\cn{Institution}}
\begin{lstlisting}
:Researcher rdfs:subClassOf [
  a owl:Restriction ; owl:onProperty :worksFor ;
  owl:someValuesFrom :Institution
] .
\end{lstlisting}

\subsection*{\dl{S} --- shorthand for $\ALC$ with transitive roles}
Not a constructor but an abbreviation for $\ALC_{R^+}$. Nearly every
practically relevant DL name starts with it.
\dlaxiom{\rn{hasAncestor} \circ \rn{hasAncestor} \sub \rn{hasAncestor}}
\begin{lstlisting}
:hasAncestor a owl:TransitiveProperty .
\end{lstlisting}

\subsection*{\dl{H} --- role hierarchy}
Subsumption between roles.
\dlaxiom{\rn{hasMother} \sub \rn{hasParent}}
\begin{lstlisting}
:hasMother rdfs:subPropertyOf :hasParent .
\end{lstlisting}

\subsection*{\dl{R} --- complex role inclusions (subsumes \dl{H})}
General role chains, subject to a regularity (acyclicity) condition, plus
reflexive, irreflexive and disjoint roles and local reflexivity
$\exists R.\mathsf{Self}$.
\begin{align*}
  \rn{hasParent} \circ \rn{hasParent} &\sub \rn{hasGrandparent}\\
  \mathsf{Disj}(\rn{hasParent}&, \rn{hasSpouse})\\
  \cn{Narcissist} &\eqv \exists \rn{loves}.\mathsf{Self}
\end{align*}
\begin{lstlisting}
:hasGrandparent owl:propertyChainAxiom ( :hasParent :hasParent ) .
:hasParent owl:propertyDisjointWith :hasSpouse .
:hasPart   a owl:ReflexiveProperty .
:hasParent a owl:IrreflexiveProperty .

:Narcissist owl:equivalentClass [
  a owl:Restriction ; owl:onProperty :loves ;
  owl:hasSelf "true"^^xsd:boolean
] .
\end{lstlisting}

\subsection*{\dl{O} --- nominals}
Individuals used inside concept expressions: enumerations $\{a,b,c\}$ and the
singleton case $\exists R.\{a\}$.
\begin{align*}
  \cn{Continent} &\eqv \{\ind{Africa}, \ind{Americas}, \ind{Antarctica},
                        \ind{Asia}, \ind{Europe}, \ind{Oceania}\}\\
  \cn{Italian} &\eqv \exists \rn{hasNationality}.\{\ind{Italy}\}
\end{align*}
\begin{lstlisting}
:Continent owl:equivalentClass [
  a owl:Class ;
  owl:oneOf ( :Africa :Americas :Antarctica :Asia :Europe :Oceania )
] .

:Italian owl:equivalentClass [
  a owl:Restriction ; owl:onProperty :hasNationality ;
  owl:hasValue :Italy
] .
\end{lstlisting}

\subsection*{\dl{I} --- inverse roles}
$\inv{R}$ as a role expression, usable wherever a role name is.
\begin{align*}
  \rn{hasChild} &\eqv \inv{\rn{hasParent}}\\
  \cn{Child} &\sub \exists \inv{\rn{hasChild}}.\cn{Person}
\end{align*}
\begin{lstlisting}
:hasChild owl:inverseOf :hasParent .

:Child rdfs:subClassOf [
  a owl:Restriction ;
  owl:onProperty [ owl:inverseOf :hasChild ] ;
  owl:someValuesFrom :Person
] .
\end{lstlisting}

\subsection*{\dl{N} --- unqualified cardinality restrictions}
$\geq n\,R$, $\leq n\,R$, $=\,n\,R$ with no filler (implicitly $\top$).
\dlaxiom{\cn{Person} \sub {\leq}\,2\ \rn{hasParent}}
\begin{lstlisting}
:Person rdfs:subClassOf [
  a owl:Restriction ; owl:onProperty :hasParent ;
  owl:maxCardinality "2"^^xsd:nonNegativeInteger
] .
\end{lstlisting}

\subsection*{\dl{Q} --- qualified cardinality restrictions}
$\geq n\,R.C$: the filler is constrained. Strictly more expressive than \dl{N}
(which is \dl{Q} with $C = \top$). Introduced in OWL~2.
\dlaxiom{\cn{Hand} \sub {\geq}\,5\ \rn{hasPart}.\cn{Finger}}
\begin{lstlisting}
:Hand rdfs:subClassOf [
  a owl:Restriction ; owl:onProperty :hasPart ;
  owl:minQualifiedCardinality "5"^^xsd:nonNegativeInteger ;
  owl:onClass :Finger
] .
\end{lstlisting}

\subsection*{\dl{F} --- functional roles}
The special case $\top \sub {\leq}\,1\,R$. Weaker than \dl{N}; present in
$\SHIF$ (OWL~Lite) precisely to keep complexity down.
\begin{align*}
  \top &\sub {\leq}\,1\ \rn{hasBiologicalMother}\\
  \top &\sub {\leq}\,1\ \inv{\rn{hasSSN}}
\end{align*}
\begin{lstlisting}
:hasBiologicalMother a owl:FunctionalProperty .
:hasSSN a owl:InverseFunctionalProperty .
\end{lstlisting}

\subsection*{$(\mathcal{D})$ --- concrete domains / datatypes}
Datatype properties and datatype ranges, including facet restrictions.
\dlaxiom{\cn{Adult} \eqv \cn{Person} \sqcap \exists \rn{hasAge}.({\geq}18)}
\begin{lstlisting}
:Adult owl:equivalentClass [
  a owl:Class ;
  owl:intersectionOf (
    :Person
    [ a owl:Restriction ; owl:onProperty :hasAge ;
      owl:someValuesFrom [
        a rdfs:Datatype ;
        owl:onDatatype xsd:integer ;
        owl:withRestrictions ( [ xsd:minInclusive 18 ] )
      ]
    ]
  )
] .
\end{lstlisting}

%%%%%%%%%%%%%%%%%%%%%%%%%%%%%%%%%%%%%%%%%%%%%%%%%%%%%%%%%%%%%%%%%%%%%%%%%%%%%%%
\section{The \dlliteplain{} family}
\label{sec:dllite}

\dlliteplain{} does not fit the letter-concatenation scheme of the previous
sections, and that is deliberate. The families above were obtained
\emph{top-down}, by starting from a rich logic and asking how much can be
retained while decidability holds. \dlliteplain{} was designed \emph{bottom-up}
(Calvanese et al., \emph{JAR} 2007) against a single criterion: conjunctive
query answering must remain \emph{first-order rewritable}, i.e.\ reducible to
the evaluation of an SQL query over the data alone. Every restriction below
follows from that criterion rather than from a menu of constructors.

\subsection*{Syntax}

The language is built from \emph{basic} concepts and roles:
\begin{align*}
  B &::= A \;\mid\; \exists R          &&\text{basic concept}\\
  R &::= P \;\mid\; \inv{P}            &&\text{basic role}\\
  C &::= B \;\mid\; \neg B             &&\text{general concept}\\
  E &::= R \;\mid\; \neg R             &&\text{general role}
\end{align*}
Here $\exists R$ is \emph{unqualified}: it abbreviates $\exists R.\top$.
A TBox contains inclusions $B \sub C$ (and, in the extensions below,
$R \sub E$); an ABox contains assertions $A(a)$ and $P(a,b)$.

The asymmetry is the essential point. Only a \emph{basic} concept may occur on
the left of $\sub$, and negation may occur only on the right. This guarantees
that the ontology has a canonical model that can be folded into the query
instead of being materialised over the data.

\subsection*{$\dllite{\mathsf{core}}$ --- concept inclusions and disjointness}

The common core: inclusions between basic concepts, plus disjointness.
Domain and range assertions are not primitives but derived readings of
$\exists P \sub A$ and $\exists \inv{P} \sub A$.
\begin{align*}
  \cn{Professor} &\sub \exists \rn{teaches}
    &&\text{every professor teaches something}\\
  \exists \rn{teaches} &\sub \cn{Professor}
    &&\text{domain of } \rn{teaches}\\
  \exists \inv{\rn{teaches}} &\sub \cn{Course}
    &&\text{range of } \rn{teaches}\\
  \cn{Professor} &\sub \neg\cn{Student}
    &&\text{disjointness}
\end{align*}
\begin{lstlisting}
:Professor rdfs:subClassOf [
  a owl:Restriction ; owl:onProperty :teaches ;
  owl:someValuesFrom owl:Thing
] .

:teaches rdfs:domain :Professor ;
         rdfs:range  :Course .

:Professor owl:disjointWith :Student .
\end{lstlisting}

\noindent
Strictly, $\dllite{\mathcal{R}}$ has only the unqualified $\exists R$. OWL~2~QL
nevertheless admits a qualified existential $\exists R.A$ in superclass
position, because it is a conservative extension: $A \sub \exists R.B$ can be
rewritten with a fresh role as
$A \sub \exists R_B$, $R_B \sub R$, $\exists \inv{R_B} \sub B$.

\subsection*{$\dllite{\mathcal{F}}$ --- functionality}

Adds functionality assertions $(\mathsf{funct}\ R)$ on basic roles, inverses
included. No role inclusions.
\begin{align*}
  \top &\sub {\leq}\,1\ \rn{hasSupervisor}
    &&(\mathsf{funct}\ \rn{hasSupervisor})\\
  \top &\sub {\leq}\,1\ \inv{\rn{hasMatricula}}
    &&(\mathsf{funct}\ \inv{\rn{hasMatricula}})
\end{align*}
\begin{lstlisting}
:hasSupervisor a owl:FunctionalProperty .
:hasMatricula  a owl:InverseFunctionalProperty .
\end{lstlisting}

\subsection*{$\dllite{\mathcal{R}}$ --- role inclusions}

Adds inclusions between roles, $R \sub E$, hence both role hierarchies and role
disjointness. No functionality. This is the logic underlying OWL~2~QL.
\begin{align*}
  \rn{hasSupervisor} &\sub \rn{knows}\\
  \inv{\rn{hasSupervisor}} &\sub \rn{supervises}\\
  \rn{hasSupervisor} &\sub \neg\rn{hasStudent}
\end{align*}
\begin{lstlisting}
:hasSupervisor rdfs:subPropertyOf :knows ;
               owl:propertyDisjointWith :hasStudent .
:supervises owl:inverseOf :hasSupervisor .
\end{lstlisting}

\subsection*{$\dllite{\mathcal{A}}$ --- combining the two, safely}

Putting $\mathcal{F}$ and $\mathcal{R}$ together without care destroys the
guarantee: with unrestricted interaction between functionality and role
inclusions, conjunctive query answering is no longer FOL-rewritable and
reasoning becomes ExpTime-hard (Artale et al., \emph{JAIR} 2009).
$\dllite{\mathcal{A}}$ restores rewritability with a syntactic proviso:
\begin{quote}
  a functional role may never be specialised --- if $(\mathsf{funct}\ R)$ or
  $(\mathsf{funct}\ \inv{R})$, then $R$ and $\inv{R}$ must not occur on the
  right-hand side of any role inclusion.
\end{quote}
$\dllite{\mathcal{A}}$ also separates abstract objects from data values,
distinguishing concepts and roles from \emph{attributes} and
\emph{value-domains} --- the direct ancestor of OWL's split between object and
datatype properties.

This explains a design choice that otherwise looks arbitrary: OWL~2~QL keeps
role inclusions and drops \texttt{owl:Functional\-Property} and
\texttt{owl:Inverse\-Functional\-Property} altogether, rather than attempting to
police the interaction at the profile level.

\subsection*{The Horn and Boolean extensions}

Artale et al.\ generalise the family along a second axis, by relaxing what may
appear on the left of $\sub$ and which Boolean operators are available. The
fragments retained for OWL~2~QL are the \emph{core} ones; the others show
precisely what each relaxation costs.

\begin{table}[h]
\centering
\footnotesize
\begin{tabular}{@{}lllp{4.1cm}@{}}
\toprule
\textbf{Fragment} & \textbf{Satisfiability} & \textbf{CQ data compl.} &
\textbf{Left-hand side of $\sub$}\\
\midrule
$\dllite{\mathsf{core}}$, $\dllite{\mathcal{F}}$, $\dllite{\mathcal{R}}$
  & NLogSpace-complete & $AC^0$ & a single basic concept $B$\\
$\dllite{\mathsf{krom}}$
  & NLogSpace-complete & $AC^0$ & binary clauses ($B_1 \sub \neg B_2$ etc.)\\
$\dllite{\mathsf{horn}}$
  & PTime-complete & PTime & conjunction $B_1 \sqcap \dots \sqcap B_n$\\
$\dllite{\mathsf{bool}}$
  & NP-complete & coNP & arbitrary Boolean combination\\
\bottomrule
\end{tabular}
\end{table}

\noindent
The jump from $AC^0$ to PTime at the Horn line is the one that matters in
practice: it is exactly the point at which the TBox can no longer be compiled
away into the query, and a reasoner must instead touch the data.

\subsection*{Why FOL-rewritability is the point}

Given a conjunctive query $q$ and a TBox $\mathcal{T}$, the \textsc{PerfectRef}
algorithm rewrites $q$ into a union of conjunctive queries $q_{\mathcal{T}}$
such that, for every ABox $\mathcal{A}$,
\[
  \mathsf{cert}(q, \langle \mathcal{T}, \mathcal{A}\rangle)
  \;=\;
  \mathsf{ans}(q_{\mathcal{T}}, \mathcal{A}),
\]
where the right-hand side evaluates $q_{\mathcal{T}}$ over $\mathcal{A}$ read as
a plain relational database, with no reasoning at all. The ontology is compiled
into the query, never into the data; the data may stay in an untouched legacy
RDBMS behind a set of mappings. This is the basis of OBDA systems such as
Mastro and Ontop, and the reason OWL~2~QL exists as a profile.

The cost is paid elsewhere: the rewriting may grow exponentially in the size of
the \emph{query} (though not of the data), which is why practical systems invest
heavily in optimising and pruning $q_{\mathcal{T}}$.

%%%%%%%%%%%%%%%%%%%%%%%%%%%%%%%%%%%%%%%%%%%%%%%%%%%%%%%%%%%%%%%%%%%%%%%%%%%%%%%
\section{Named DLs and their OWL counterparts}

\begin{table}[h]
\centering
\begin{tabular}{@{}lll@{}}
\toprule
\textbf{Description Logic} & \textbf{OWL language} &
\textbf{Complexity (concept sat.)}\\
\midrule
$\EL$ / $\EL^{++}$          & OWL~2~EL    & PTime-complete\\
\dllite{\mathcal{R}}        & OWL~2~QL    & NLogSpace; $AC^0$ data compl.\ for CQs\\
\dl{DLP} / $pD^*$-style     & OWL~2~RL    & PTime-complete\\
$\ALC$                      & ---         & ExpTime-complete\\
$\SHIF^{(\mathcal{D})}$     & OWL~1~Lite  & ExpTime-complete\\
$\SHOIN^{(\mathcal{D})}$    & OWL~1~DL    & NExpTime-complete\\
$\SROIQ^{(\mathcal{D})}$    & OWL~2~DL    & N2ExpTime-complete\\
unrestricted, RDF semantics & OWL~2~Full  & undecidable\\
\bottomrule
\end{tabular}
\end{table}

\noindent
Note that OWL~Lite, despite its name, is $\SHIF^{(\mathcal{D})}$: it has the same
worst-case complexity as $\SHIF$, and most of OWL~DL can be encoded into it.
This is why it was effectively abandoned in OWL~2.

%%%%%%%%%%%%%%%%%%%%%%%%%%%%%%%%%%%%%%%%%%%%%%%%%%%%%%%%%%%%%%%%%%%%%%%%%%%%%%%
\section{The three OWL 2 profiles}

The profiles are \emph{pairwise incomparable} --- they do not form a chain.
Each is a syntactic subset of OWL~2~DL designed to make a \emph{different}
reasoning task tractable.

\paragraph{OWL 2 EL ($\EL^{++}$).}
Target: very large terminologies (SNOMED~CT, GO, NCIt).
Retains $\sqcap$, $\exists R.C$, $\top$, $\bot$, nominals, role chains,
$\mathsf{Self}$, domain and range, keys, datatypes.
Drops $\forall R.C$, $\sqcup$, $\neg$, inverse roles, all cardinality
restrictions, functional properties, and disjunctive reasoning.
Classification is polynomial.

\paragraph{OWL 2 QL (\dllite{\mathcal{R}}, see Section~\ref{sec:dllite}).}
Target: query answering over large ABoxes stored in relational databases via
rewriting (OBDA, Ontop). The syntax is asymmetric: on the left of $\sub$ one may
use a class name or $\exists R.\top$; on the right, a class name,
$\exists R.C$, or a negation. Drops $\sqcup$, $\exists R.C$ in subclass
position, property chains, functional and transitive properties, nominals, and
keys. Conjunctive queries are FOL-rewritable into SQL, giving $AC^0$ data
complexity.

\paragraph{OWL 2 RL.}
Target: rule engines and forward chaining over RDF (partial axiomatisation as
Datalog). Also asymmetric: subclass position allows $\sqcap$, $\sqcup$,
$\exists R.C$ and nominals; superclass position allows $\sqcap$,
$\forall R.C$, ${\leq}1\,R.C$, ${\leq}0\,R.C$ and negation. Drops
$\exists R.C$ on the right (no existential blank-node generation) and $\sqcup$
on the right. Keeps inverses, transitivity, functionality, property chains and
keys. PTime data complexity, and it degrades gracefully: the RL rule set can be
run over any RDF graph, including OWL~Full ones, at the cost of incompleteness.

\begin{table}[h]
\centering
\footnotesize
\begin{tabular}{@{}lccc@{}}
\toprule
\textbf{Construct} & \textbf{EL} & \textbf{QL} & \textbf{RL}\\
\midrule
$C \sqcap D$                       & yes & partial & yes\\
$C \sqcup D$                       & no  & no      & subclass position only\\
$\neg C$                           & no  & superclass position only & superclass position only\\
$\exists R.C$                      & yes & superclass position only & subclass position only\\
$\forall R.C$                      & no  & no      & superclass position only\\
Nominals $\{a\}$                   & yes & no      & yes\\
Inverse roles $\inv{R}$            & no  & yes     & yes\\
Cardinality restrictions           & no  & no      & ${\leq}1$, ${\leq}0$ only\\
Functional properties              & no  & no      & yes\\
Transitive properties              & yes & no      & yes\\
Property chains                    & yes & no      & yes\\
$\exists R.\mathsf{Self}$          & yes & no      & no\\
Keys                               & yes & no      & yes\\
\midrule
Data complexity                    & PTime & $AC^0$ & PTime\\
\bottomrule
\end{tabular}
\end{table}

%%%%%%%%%%%%%%%%%%%%%%%%%%%%%%%%%%%%%%%%%%%%%%%%%%%%%%%%%%%%%%%%%%%%%%%%%%%%%%%
\section{Rule-based entailment in triple stores: the GraphDB rulesets}
\label{sec:graphdb}

Triple stores rarely embed a DL reasoner. GraphDB, like most of them, offers
entailment through \emph{rulesets}: \texttt{empty}, \texttt{rdfs},
\texttt{rdfsplus}, \texttt{owl-horst}, \texttt{owl-max}, \texttt{owl2-ql} and
\texttt{owl2-rl}, each also available as an \texttt{-optimized} variant.
The names suggest a correspondence with the languages of the previous
sections, but the correspondence is loose, for three reasons.

\begin{itemize}
  \item \textbf{The engine is a rule engine.} Each ruleset is a set of Horn
    rules (in GraphDB's \texttt{.pie} format) applied by forward chaining; the
    closure is fully materialised at load time.
  \item \textbf{The semantics is RDF-based.} Rules match triples, not axioms;
    classes may be instances, and no Direct-Semantics interpretation is
    involved.
  \item \textbf{There is no syntactic check.} A ruleset is applied to any
    graph, whether or not it lies in the corresponding profile or DL.
\end{itemize}

\noindent
The consequence is that none of the rulesets implements a full description
logic. Each is best understood as a \emph{Horn fragment} of some DL: sound, but
incomplete with respect to the full logic whose vocabulary it accepts.

\subsection*{Mapping to DL families}

\begin{table}[h]
\centering
\footnotesize
\begin{tabular}{@{}llp{6.4cm}@{}}
\toprule
\textbf{GraphDB ruleset} & \textbf{Closest DL / profile} &
\textbf{DL axioms covered}\\
\midrule
\texttt{rdfs}
  & $\rho$df $\subset \dllite{\mathcal{R}}$
  & $A \sub B$, $P \sub Q$, $\exists P \sub A$ (domain),
    $\exists \inv{P} \sub A$ (range)\\
\texttt{rdfsplus}
  & Horn fragment of \dl{SHI} ($\approx$ DLP)
  & $+$ $P \eqv \inv{Q}$, $P \eqv \inv{P}$ (symmetry), $\mathsf{Trans}(P)$\\
\texttt{owl-horst}
  & $pD^*$ $\approx$ Horn-\dl{SHOIF}
  & $+$ $a \approx b$, $\mathsf{Func}(P)$, $\mathsf{Func}(\inv{P})$,
    $\exists P.C \sub A$, $A \sub \forall P.C$, $A \eqv \exists P.\{o\}$,
    $A \eqv B$\\
\texttt{owl-max}
  & Horn fragment of $\SHIF^{(\mathcal{D})}$ (OWL~Lite syntax)
  & $+$ $B \sqcup C \sub A$, intersections, ${\leq}1\,P$\\
\texttt{owl2-rl}
  & OWL~2~RL $\subset$ Horn-$\SROIQ^{(\mathcal{D})}$
  & $+$ property chains, keys, disjoint properties, qualified ${\leq}1$\\
\texttt{owl2-ql}
  & OWL~2~QL $\approx \dllite{\mathcal{R}}$
  & $A \sub B$, $A \sub \exists P.B$, role inclusions, inverses,
    disjointness\\
\bottomrule
\end{tabular}
\end{table}

\noindent
The reference points are $\SHIF^{(\mathcal{D})}$ for OWL~Lite,
$\SHOIN^{(\mathcal{D})}$ for OWL~1~DL and $\SROIQ^{(\mathcal{D})}$ for OWL~2~DL.
GraphDB ships no ruleset for OWL~2~EL.

\subsection*{Why every ruleset is a Horn fragment}

Two limitations are shared by all rulesets, and together they account for most
of the distance from the full logics.

\begin{itemize}
  \item \textbf{No disjunctive consequences.} A rule can derive $A(a)$ but never
    $A(a) \lor B(a)$. For a union
    $\cn{Staff} \eqv \cn{Academic} \sqcup \cn{TechnicalStaff}$ only the
    direction $\cn{Academic} \sub \cn{Staff}$ is usable; the converse, which a
    tableau would exploit by case analysis, is lost.
  \item \textbf{No new individuals.} Rules never invent a witness. An axiom
    $A \sub \exists P.B$ is accepted and stored, but it yields no
    $P$-successor, and nothing that would follow from one.
\end{itemize}

\subsection*{Ruleset by ruleset}

\paragraph{RDFS.} Not a DL at all, since it permits metamodelling. Its
DL-expressible part is a small subset of both $\dllite{\mathcal{R}}$ and \EL{}.

\paragraph{RDFS-Plus.} Adds the role axioms of \dl{SHI} --- inverse and
transitive roles on top of the role hierarchy --- but no new concept
constructors. It lies within Description Logic Programs, the intersection of
DL and Horn logic identified by Grosof et al.\ (WWW 2003).

\paragraph{OWL-Horst.} Ter Horst's $pD^*$ (\emph{JWS} 2005) was defined as a
semantic extension of RDFS with an ``if''-semantics for the OWL vocabulary, not
as a DL. Its constructors roughly cover Horn-\dl{SHOIF}: \texttt{owl:hasValue}
contributes nominals, functional and inverse-functional properties contribute
\dl{F}. It is sound but incomplete for \SHOIN{}. GraphDB's implementation goes
beyond textbook $pD^*$ in handling \texttt{owl:intersectionOf} in both
directions (see the empirical check below).

\paragraph{OWL-Max.} Accepts OWL~Lite, i.e.\ $\SHIF^{(\mathcal{D})}$, syntax, but
derives only its Horn-expressible consequences. Relative to OWL-Horst it adds,
chiefly, the subclass direction of \texttt{owl:unionOf}, rule-expressible
cardinality (above all ${\leq}1$, which behaves as a local functionality
constraint and so produces \texttt{owl:sameAs}), and further consistency
checks. The exact rule difference is version-dependent; the authoritative
source is a diff of the two \texttt{.pie} files.

\paragraph{OWL 2 RL.} The W3C profile was itself derived from DLP and $pD^*$.
On ontologies that satisfy the RL syntactic restrictions, the RL/RDF rules are
sound and complete for ground atomic consequences (Theorem~PR1 of the OWL~2
Profiles specification); outside them they remain sound but incomplete. Data
complexity is polynomial.

\paragraph{OWL 2 QL.} The $\dllite{\mathcal{R}}$ profile of
Section~\ref{sec:dllite}, designed for query rewriting with $AC^0$ data
complexity. GraphDB does \emph{not} rewrite: it materialises the QL
consequences instead, which forgoes the main benefit of the profile and, as the
empirical check shows, produces far more implicit triples than the other
rulesets. The ruleset is also more permissive than the profile.

\subsection*{The \texttt{-optimized} variants}

According to the GraphDB documentation, the optimised rulesets remove
inferences that matter for formal completeness but are rarely queried: the
axiomatic triples and rule conclusions involving \texttt{rdfs:Resource}, and
a handful of further axiomatic triples (e.g.\ \texttt{owl:sameAs} typed as
symmetric and transitive). The documented list does \emph{not} include domain
and range rules --- which makes the first finding below unexpected.

\subsection*{Empirical check}

The claims above were tested on a small university knowledge graph that
exercises one construct per axiom block, loaded into four repositories
(\texttt{rdfsplus-optimized}, \texttt{owl-horst-optimized},
\texttt{owl-max-optimized}, \texttt{owl2-ql-optimized}; \texttt{owl:sameAs}
enabled) and queried with inference switched on. The key trick is an
individual with no asserted type, whose every class membership must therefore
be inferred:
\begin{lstlisting}
ex:teaches a owl:ObjectProperty ;
           rdfs:domain ex:Academic ; rdfs:range ex:Course .
ex:alice   ex:teaches ex:kg101 .     # no rdf:type for ex:alice
\end{lstlisting}

\begin{table}[h]
\centering
\footnotesize
\newcommand{\tyes}{$\checkmark$}
\newcommand{\tno}{$\times$}
\begin{tabular}{@{}llcccc@{}}
\toprule
& \textbf{Construct tested} & \textbf{RDFS-Plus} & \textbf{Horst} &
  \textbf{Max} & \textbf{QL}\\
\midrule
Q1 & domain, subclass                        & \tno$^{a}$ & \tyes & \tyes & \tyes\\
Q2 & subproperty, range                      & \tno$^{a}$ & \tyes & \tyes & \tyes\\
Q3 & inverse, symmetric                      & \tyes & \tyes & \tyes & \tyes\\
Q3 & transitive                              & \tyes & \tyes & \tyes & \tno\\
Q4 & equivalent class / property            & --$^{b}$ & \tyes & \tyes & \tyes\\
Q5 & sameAs via (inverse-)functional prop.   & \tno & \tyes & \tyes & \tno\\
Q6 & hasValue, someValuesFrom, allValuesFrom & \tno & \tyes & \tyes & \tno\\
Q7 & intersection, subclass side             & \tno & \tyes$^{c}$ & \tyes & \tyes$^{d}$\\
Q7 & intersection, superclass side           & \tno & \tyes & \tyes & \tyes\\
Q7 & union, subclass side                    & \tno & \tno & \tyes & \tno\\
Q8 & property chain                          & \tno & \tno & \tno & \tno\\
Q9 & keys                                    & \tno & \tno & \tno & \tno\\
\midrule
Q10 & implicit triples                      & 193 & 337 & 361 & 1596\\
\bottomrule
\end{tabular}

\smallskip
\raggedright
$^{a}$~unexpected, see below.\quad
$^{b}$~not observable: the query depends on the domain inference.\\
$^{c}$~beyond textbook $pD^*$.\quad
$^{d}$~beyond the OWL~2~QL profile.
\end{table}

\noindent
Most results match the analysis. The OWL-Horst and OWL-Max columns agree
everywhere except on union, which confirms both the inclusion between the two
and that union is the only difference this test bed exposes. Property values
stated on an alias propagate across \texttt{owl:sameAs}. OWL~2~QL has no
transitivity, no equality and no \texttt{hasValue}, and does not derive the
qualified existential on the subclass side. Q8 and Q9 are empty everywhere
simply because no \texttt{owl2-rl} repository was part of the run; that
ruleset should return the chained affiliations and the key-based equality.

Four observations deviate from expectations.
\begin{enumerate}
  \item \textbf{\texttt{rdfsplus-optimized} applies neither domain nor
    range.} Inverse, symmetric and transitive properties work, and
    equivalent classes are turned into subclass links, so the ruleset is
    active; yet \texttt{ex:alice} receives no type and \texttt{ex:kg101} is
    never a course. An earlier run on an RDFS repository showed the same
    symptom. Nothing in the documentation of the optimised variants accounts
    for it. To be confirmed by querying range and subproperty inferences
    separately, and by repeating the run on the non-optimised ruleset.
  \item \textbf{OWL-Horst supports intersection in both directions}, which
    textbook $pD^*$ does not.
  \item \textbf{OWL~2~QL derives membership in an intersection} (an individual
    typed with both conjuncts becomes an instance of the intersection class),
    a subclass-side use that $\dllite{\mathcal{R}}$ does not allow.
  \item \textbf{OWL~2~QL materialises almost five times as much as OWL-Max},
    partly because it types every individual as \texttt{owl:Thing}, which the
    other optimised rulesets do not.
\end{enumerate}

\subsection*{Practical consequences}

\begin{itemize}
  \item \textbf{Names are not guarantees.} A ruleset may be weaker than its
    name suggests (RDFS-Plus above) or stronger than the profile it is named
    after (OWL-Horst, OWL~2~QL). Test the constructs an application relies on,
    rather than inferring support from the ruleset name.
  \item \textbf{DL-illegal input is silently accepted.} An
    \texttt{owl:InverseFunctionalProperty} on a datatype property is outside
    \SHOIN{} and \SROIQ{}, yet it produces \texttt{owl:sameAs} in every ruleset
    that handles inverse functionality.
  \item \textbf{The gap to a DL reasoner is measurable.} Running the same
    ontology through HermiT or Konclude (or ELK for its EL part) and comparing
    inferred class assertions isolates exactly the consequences that need
    disjunctive or existential reasoning.
  \item \textbf{Materialisation cost varies sharply between rulesets}, and is
    not monotone in apparent expressivity, as the QL figure shows.
\end{itemize}

%%%%%%%%%%%%%%%%%%%%%%%%%%%%%%%%%%%%%%%%%%%%%%%%%%%%%%%%%%%%%%%%%%%%%%%%%%%%%%%
\section{Openly available reasoners}
\label{sec:reasoners}

The sections above describe languages; this one describes what is actually
implemented. Three warnings before the tables.

\emph{First, ``supports \dl{X}'' is a claim about the calculus, not about the
code.} Most systems implement a strict subset of the logic they advertise, and
the gap is usually in datatypes, nominals or the more exotic role axioms. The
\emph{Notes} column flags the cases where the gap is large enough to change a
tool choice.

\emph{Second, expressivity is rarely the binding constraint in practice.}
Maintenance is. A systematic 2023 review of the field found that fewer than half
of the analysed reasoners had seen any development activity in the preceding
three years, and that several of the best-known names --- HermiT among them ---
were effectively abandoned while remaining in daily use for want of
alternatives (Abicht, arXiv:2309.06888).

\emph{Third, the comparative evidence is dated.} The last edition of the OWL
Reasoner Evaluation competition was ORE~2015 (Parsia et al., \emph{JAR} 2017);
there has been no comparable successor, so performance claims newer than that
rest on individual papers rather than on a common benchmark.

\subsection*{OWL 2 DL: the $\SROIQ$ tier}

\begin{table}[h]
\centering
\footnotesize
\begin{tabular}{@{}lllp{4.5cm}@{}}
\toprule
\textbf{Reasoner} & \textbf{Expressivity} & \textbf{Method} & \textbf{Notes}\\
\midrule
Konclude & $\SROIQ\mathcal{V}^{(\mathcal{D})}$ & tableau + saturation
  & C++/Qt, LGPLv3. $\mathcal{V}$ = nominal schemas, i.e.\ DL-safe rules in
    ontology syntax. Won most DL tracks of ORE~2015\\
HermiT & $\SROIQ^{(\mathcal{D})}$ & hypertableau
  & Java, LGPL. Upstream dead since 1.3.8; the living artefact is the OWL~API
    fork, explicitly unsupported by its authors. Ships with Prot\'eg\'e\\
FaCT++ & $\SROIQ^{(\mathcal{D})}$ & optimised tableau
  & C++, LGPL. Last commit 2017; still ships with Prot\'eg\'e\\
JFact & $\SROIQ^{(\mathcal{D})}$ & tableau
  & Java port of FaCT++, maintained under the OWL~API umbrella\\
Openllet & $\SROIQ^{(\mathcal{D})}$ & tableau
  & Java. Community fork of Pellet; sporadic but real updates\\
Pellet & $\SROIQ^{(\mathcal{D})}$ & tableau
  & Java. Historically important, abandoned; use Openllet instead\\
Sequoia & $\dl{SRIQ}$ & consequence-based
  & Scala, GPL-3. No nominals --- the trade for consequence-based
    scalability at high expressivity\\
Racer & $\dl{SRIQ}^{(\mathcal{D})}$ & tableau
  & Common Lisp, BSD-3. Source on GitHub, project abandoned\\
Vampire & FOL & superposition
  & C++. A general theorem prover; OWL~DL reached by translation, not a
    native OWL tool\\
\bottomrule
\end{tabular}
\end{table}

\noindent
For new work at this tier the realistic choices are Konclude, which is the
fastest and the only one under active development, and HermiT, which is native
to the OWL API and ubiquitous but unmaintained. Everything else is either a
fork of one of these or of historical interest only.

\subsection*{OWL 2 EL: the $\EL^{++}$ tier}

\begin{table}[h]
\centering
\footnotesize
\begin{tabular}{@{}lllp{4.5cm}@{}}
\toprule
\textbf{Reasoner} & \textbf{Expressivity} & \textbf{Method} & \textbf{Notes}\\
\midrule
ELK & $\EL^{++}$, safe nominals & concurrent consequence-based
  & Java, Apache~2.0. Partial support only for datatypes, $\neg$ and $\sqcup$;
    incremental reclassification. The de-facto standard for biomedical
    terminologies\\
jcel & $\EL^{+}$ & completion rules
  & Java. Subset of OWL~2~EL, Prot\'eg\'e plugin\\
CEL & $\EL^{+}$ & completion rules
  & The original large-scale EL reasoner; now of mainly historical interest\\
ElepHant & $\EL^{+}$ & consequence-based
  & C. Built for small memory and CPU budgets, e.g.\ embedded use\\
Whelk & $\EL$ + RL & consequence-based
  & Scala. Designed for large EL TBoxes combined with RL-shaped ABoxes;
    integrated with ROBOT in the OBO toolchain\\
Snorocket & OWL~2~EL subset & completion rules
  & Java. Adds concrete domains; developed around SNOMED~CT tooling\\
\bottomrule
\end{tabular}
\end{table}

\noindent
ELK's own documentation is unusually honest about its boundaries: the supported
fragment is $\EL^{++}$ restricted to \emph{safe} nominals, with only partial
support for concrete roles, \texttt{Object\-Complement\-Of} and
\texttt{Object\-Union\-Of}, and it warns about unsupported constructs at load time
rather than failing silently. An ontology that is valid OWL~2~EL on paper may
therefore still be classified incompletely --- worth checking before attributing
a missing inference to a modelling error.

\subsection*{OWL 2 QL: query rewriting and OBDA}

\begin{table}[h]
\centering
\footnotesize
\begin{tabular}{@{}lllp{4.5cm}@{}}
\toprule
\textbf{Reasoner} & \textbf{Expressivity} & \textbf{Method} & \textbf{Notes}\\
\midrule
Ontop & $\dllite{\mathcal{R}}$ + RDFS & rewriting to SQL
  & Java, Apache~2.0, actively developed (5.5.0, February 2026). SPARQL over
    virtual knowledge graphs via R2RML mappings; the reference OBDA system\\
Mastro & $\dllite{\mathcal{A}}$ & rewriting to SQL
  & Sapienza. Academic licence, not open source\\
Clipper & Horn-$\dl{SHIQ}$ & rewriting to Datalog
  & Java. Above QL in expressivity but still rewritable; abandoned, and its
    Datalog back end (DLV) is no longer distributed\\
Requiem & $\dl{ELHIO}$ & resolution-based rewriting
  & Prototype accompanying the rewriting algorithm; not maintained\\
\bottomrule
\end{tabular}
\end{table}

\noindent
This tier is effectively a single-vendor landscape: Ontop is the only openly
available, maintained implementation of the rewriting approach described in
Section~\ref{sec:dllite}.

\subsection*{OWL 2 RL: rules and materialisation}

\begin{table}[h]
\centering
\footnotesize
\begin{tabular}{@{}lllp{4.5cm}@{}}
\toprule
\textbf{Reasoner} & \textbf{Expressivity} & \textbf{Method} & \textbf{Notes}\\
\midrule
RDFox & OWL~2~RL, Datalog, SWRL & parallel materialisation
  & C++. Highest performance in this class, but commercial and closed source;
    free academic licences available\\
VLog & Datalog + existential rules & chase
  & C++/Java, open source. Rule engine rather than an OWL tool; OWL~RL reached
    by translating axioms to rules\\
OWL-RL & OWL~2~RL & forward chaining
  & Python, on top of RDFLib. Convenient rather than fast\\
reasonable & OWL~2~RL subset & forward chaining
  & Rust, with Python bindings. Substantially faster than OWL-RL\\
Arachne & OWL~RL & Rete
  & Scala. Built for Gene Ontology causal activity models\\
Apache Jena & RDFS, partial OWL & forward/backward rules
  & Java. Deliberately incomplete rule sets; ubiquitous, and often what people
    mean when they say ``the reasoner in my triple store''\\
GraphDB & RDFS to OWL~2~RL / QL & forward chaining
  & Java. Commercial, with a free edition. Rulesets are Horn fragments and
    deviate from their names in both directions; see
    Section~\ref{sec:graphdb}\\
EYE & N3 rules & Euler-path chaining
  & Prolog/WebAssembly, very actively developed. Not an OWL reasoner as such,
    but widely used for OWL-adjacent rule reasoning\\
\bottomrule
\end{tabular}
\end{table}

\subsection*{Hybrid and pay-as-you-go systems}

A fourth group does not fit the profile tiers: these systems combine reasoners
rather than implementing a calculus.

\begin{itemize}
  \item \textbf{PAGOdA} answers conjunctive queries over OWL~2~DL by handling
    the RL-shaped bulk with a Datalog engine (RDFox) and delegating only the
    residual hard part to a fully-fledged DL reasoner (HermiT). Complete where
    the data allows, and fast where it matters.
  \item \textbf{MORe} and \textbf{Chainsaw} split an ontology into modules and
    route each to the cheapest adequate reasoner --- typically ELK for the EL
    part and HermiT for the remainder. Both are now unmaintained.
  \item \textbf{TrOWL} approximates an OWL~2~DL ontology into QL or EL,
    trading completeness for tractability in a principled way.
\end{itemize}

\subsection*{Choosing one}

Expressivity narrows the field; the remaining criteria usually decide it.
\begin{itemize}
  \item \textbf{Large terminology, EL-shaped} --- ELK. Nothing else is
    competitive on SNOMED-scale classification, and the incremental mode
    changes the editing workflow.
  \item \textbf{Full OWL~2~DL, performance matters} --- Konclude. If the
    deployment must be pure Java and OWL~API-native, HermiT, with the
    maintenance caveat understood.
  \item \textbf{Data in an RDBMS, ontology used as a query interface} ---
    Ontop.
  \item \textbf{Large RDF graphs, forward chaining acceptable} --- RDFox if the
    licence is acceptable, otherwise a rule engine (VLog, Arachne) or the
    materialisation built into the triple store.
  \item \textbf{Embedded or constrained hardware} --- ElepHant (EL) or LiRoT
    (RL subset); both were built explicitly for that envelope.
\end{itemize}

\end{document}
